% Transcription — fonds Alexandre Grothendieck, shelfmark 14, batch 1 (pages 1-20).
% Established from the facsimile archives/batches/14/batch-01.pdf (a mirror of
% https://grothendieck.umontpellier.fr/14.pdf), per the skill
% transcribe-grothendieck. The batch file carries no cover sheet: PDF page N is
% archivists' page N; the pencilled numbers bottom-left ("1", "2", "4" ... "18",
% "20") read at the PDF page of the same number.
%
% Pass: Opus 5.5 (claude-opus-5-5), 2026-09-23 — first pass, unchecked against the pages by a human.
%
% What the batch is. Two short runs in his hand, each opened by a cover sheet
% he inscribed, written on the rectos of leaves whose versos are typescripts
% unrelated to them (an English typescript on the Brauer group of a ringed
% space, pages III.4-IV.2, at pp. 3, 5, 7, 9, 11, 13, 15, scanned upside down;
% a French typescript in blue, p. 19). Those versos carry no hand and are
% absent, as is the blank p. 17. Nothing on the leaves is dated; no letter in
% this batch.
%   p. 1    cover: « Structure d'effectivité » (his title, underlined), and a
%           boxed title struck through by diagonals (« Les éléments de
%           structure motiviques et leur interprétation en termes de groupes
%           algébriques sur Q »). No \page; his words are the section head.
%   pp. 2-10  the run itself, headed « Structure » on p. 2 and opening at item
%           « 4°) » — items 1°-3° are not in this batch (p. 8 and p. 10 refer
%           back to « 3° », the Hodge/polarisation item; « structures 3° et
%           4° »). An effectivity structure on the motives = a sub-category
%           M^+ of effective motives, stable by ⊗, sums, sub-objects, with
%           every E(-n) effective for n large; equivalently a cone K^+(G_C)
%           in the virtual-character ring K(G). Effective cocharacters
%           i : G_m,C → G_C, the sub-categories Pl(i), the pairing <i, ε>
%           (ε the weight character); cones Γ^+ in Γ = Hom(G_m, T_C) invariant
%           under W, the quotient Γ' = Γ/W = [Hom(G_m, G)/G](C) = S(C) with
%           π_1(Q, C) acting through a finite quotient; the map
%           I(G) → Eff(G), I*(G) → Pol(G) × Eff(G); questions (injectivity;
%           is every motivic effectivity structure given by an i?).
%   pp. 12-14 the case of a torus T over Q: M the character lattice with
%           Galois action, ε ∈ M^π, i ∈ M^∨π, <i, ε> = 2; the structure only
%           depends on T_R; he claims dim S = 1 (Ker ε compact) — the
%           attempted proof is cancelled; then K(T) = Z(M)^π, effective
%           characters as those ≥ 0 on S(R)°, the standard procedure through
%           i_1 : G_m → T_C and polynomiality at t → 0; finiteness of
%           effective characters of given weight fails for dim T ≥ 3.
%   p. 16   cover: « Équivalence d'Albanese et de Picard des cycles
%           algébriques ». No \page.
%   pp. 18, 20  a worked example (headed « E »): A an elliptic curve,
%           Y = A × P, Θ_0 = Δ − A×(0) − (0)×A, the codimension-2 cycle
%           z = (id_A × i)_*(Θ) on A × X^3 with X the projective closure of
%           a line bundle over Y; s_*(z^2) = −2 c(0,0), and replacing c by −c
%           gives −2 and 2, « des signes différents !! ». The computation is
%           complete within p. 20.
%
% Readings to note. Grothendieck's « 4°) » on p. 2 is an open 4, easily taken
% for a 7. Λ(-1) is the Tate twist, as in folder 12 batch 9. Underlined
% scheme-valued Hom, Aut, Autext are set upright without the underline (see
% the note on p. 6). The fast connecting prose often does not carry; the
% formulas mostly do.
\documentclass[11pt,a4paper]{article}
% Le préambule commun à toutes les transcriptions du fonds Grothendieck.
%
% Il tient en une idée : **l'incertitude se compose, elle ne se résout pas.**
% Une transcription qui lisse un mot douteux a détruit la seule chose pour
% laquelle on la faisait — un lecteur doit pouvoir savoir, à chaque endroit,
% ce qui était sur le feuillet et ce qui a été ajouté. Les six macros qui
% suivent sont donc l'appareil critique, et non de la décoration.
%
% Les mêmes commandes sont reconnues par `scripts/render.mjs`, qui produit la
% vue de lecture du site. Une macro ajoutée ici doit l'être aussi là-bas,
% faute de quoi le rendu échouera bruyamment — ce qui est voulu : un rendu
% silencieusement faux est pire qu'un rendu absent.

\ProvidesPackage{grothendieck}[2026/08/08 Transcriptions du fonds Grothendieck]

\usepackage{amsmath,amssymb,amsthm}
\usepackage{mathrsfs}
\usepackage{stmaryrd}
% Les diagrammes commutatifs sont partout dans ce fonds : c'est la langue
% même de la géométrie algébrique de Grothendieck, et une transcription qui
% les rendrait en prose ne transcrirait rien.
\usepackage{tikz-cd}
% Français par défaut — c'est la langue des feuillets — mais l'anglais est
% chargé aussi : la traduction et le résumé sont des documents anglais, et la
% typographie française (espaces avant les deux-points) y serait une faute.
% Ces documents basculent par \selectlanguage{english} en tête de corps.
\usepackage[english,french]{babel}
\usepackage{fontspec}
\usepackage{geometry}
\geometry{a4paper,margin=25mm,marginparwidth=28mm}
\usepackage{marginnote}
\usepackage{xcolor}
% ulem, not soul. `soul`'s \st and \ul are built on a character-by-character
% scan that XeLaTeX's font machinery breaks: the document fails with an
% undefined control sequence rather than degrading. ulem does the same two jobs
% and is Unicode-engine safe. `normalem` stops it hijacking \emph, which the
% transcriptions use for Grothendieck's own emphasis.
\usepackage[normalem]{ulem}
\usepackage{hyperref}
\hypersetup{colorlinks=true,linkcolor=black,urlcolor=black,pdfborder={0 0 0}}

% --- Métadonnées du lot -----------------------------------------------------
% Reprises telles quelles par le script de rendu, qui les affiche en tête de
% la vue de lecture. Elles fixent ce que le fichier prétend couvrir : sans
% elles, une transcription flotte sans point d'attache dans un fonds de seize
% mille feuillets.

\newcommand{\folder}[1]{\gdef\grothFolder{#1}}
\newcommand{\batch}[1]{\gdef\grothBatch{#1}}
\newcommand{\leaves}[2]{\gdef\grothFirst{#1}\gdef\grothLast{#2}}
\newcommand{\foldertitle}[1]{\gdef\grothFoldertitle{#1}}
\gdef\grothFolder{?}\gdef\grothBatch{?}\gdef\grothFirst{?}\gdef\grothLast{?}\gdef\grothFoldertitle{}

% --- Le résumé --------------------------------------------------------------
%
% Il ouvre la lecture modernisée, sous le titre « Résumé », et il est composé
% en petit. Ce n'est pas un ornement : c'est le seul passage du document qui
% ne soit pas des mathématiques, et un lecteur doit pouvoir le reconnaître
% avant de l'avoir lu — puis le sauter s'il connaît déjà le sujet, ou s'y
% arrêter s'il ne le connaît pas. Le filet qui le suit marque la reprise du
% corps.

\newenvironment{resume}
  {\small\setlength{\parskip}{0.45em}}
  {\par\medskip\hrule\medskip}

% --- Mots-clés ---------------------------------------------------------------
%
% En anglais, en clôture du résumé de la lecture modernisée : le vocabulaire
% moderne sous lequel chercher ce que les feuillets construisent. Le manifeste
% du site (`npm run manifest`) les extrait pour étiqueter la cote entière —
% cette ligne est la source unique des tags, il n'y a pas de fichier à côté.
\newcommand{\keywords}[1]{\par\smallskip\noindent
  {\footnotesize\textsc{Keywords} --- \textit{#1}}\par}

% --- L'appareil critique ----------------------------------------------------

% Le numéro de feuillet, en marge. C'est l'ancre qui permet de régler un
% désaccord sur une formule en regardant *un* feuillet plutôt qu'une chemise
% de six cents. Le site s'en sert aussi pour tourner les pages du fac-similé
% quand on fait défiler la transcription.
\newcommand{\leaf}[1]{%
  \marginnote{\footnotesize\color{gray!70}\textsf{#1}}%
  \label{leaf:#1}\ignorespaces}

% Illisible : on ne devine pas. Un mot inventé qui se lit comme les autres est
% le pire résultat possible.
\newcommand{\ill}{\textcolor{red!60!black}{[\,\ldots\,]}}

% Lecture douteuse : le mot est donné, et signalé comme incertain.
\newcommand{\uncertain}[1]{\uline{#1}}

% Ajout de l'éditeur — une lettre manquante, un mot sous-entendu.
\newcommand{\add}[1]{\textcolor{blue!50!black}{[#1]}}

% Biffé par Grothendieck lui-même. Ce qu'il a rayé fait partie du document :
% une rature dit souvent où la pensée a changé de direction.
\newcommand{\struck}[1]{\sout{#1}}

% Note du transcripteur, sur l'état matériel du feuillet.
\newcommand{\note}[1]{\par\smallskip{\small\color{gray!60!black}\textit{[#1]}}\par\smallskip}

% Note marginale de Grothendieck, distincte du corps du texte.
\newcommand{\marginal}[1]{\par\smallskip\noindent\hspace{1em}%
  {\small\color{gray!60!black}#1}\par\smallskip}

% --- Titre ------------------------------------------------------------------

\AtBeginDocument{%
  \begin{center}
    {\large\bfseries Fonds Alexandre Grothendieck --- cote n\textsuperscript{o}~\grothFolder}\\[2pt]
    {\itshape\grothFoldertitle}\\[4pt]
    {\small feuillets \grothFirst--\grothLast\ (lot \grothBatch)}\\[2pt]
    {\footnotesize\color{gray!60!black}Transcription établie d'après le fac-similé numérisé
     par l'Université de Montpellier.\\
     Les lectures incertaines sont soulignées, les illisibles notées [\,\ldots\,],
     les ajouts entre crochets.}
  \end{center}
  \vspace{1em}}

\endinput


\folder{14}
\batch{1}
\pages{1}{20}
\dating{1965-[vers 1971]}
\watermark{Édition de démonstration}
\foldertitle{Théorie des cycles algébriques : conjectures de Hodge, Tate, Lefschetz, Weil : notes manuscrites (s.d.), lettre (1965).}

\begin{document}

\section{Structure d'effectivité}

\note{le titre est de sa main, souligné, sur un feuillet de couverture
(feuillet 1) laissé par ailleurs blanc ; au-dessus, dans un cadre, un titre
biffé de traits diagonaux : « Les éléments de structure motiviques et leur
interprétation en termes de groupes algébriques sur $\mathbb{Q}$ »}

\page{2}
\marginal{Structure}
\note{ce mot est en tête du feuillet, de sa main}

4°) La donnée d'une sous-catégorie strict.\ pleine de $\mathcal{M}$
(motifs effectifs), stable par $\otimes$, par sommes, par
\struck{\ill{}} sous-\struck{\ill{}}objets, telle que $\Lambda(-1)$ soit
effectif et que pour tout $E$, $\exists$ entier $n$ tel que
$E(-n) = E \otimes \Lambda(-1)^{\otimes n}$ soit « effectif ».
\note{« strict.\ pleine » et « $\Lambda(-1)$ soit effectif et que » sont
interlinéaires ; devant $\Lambda(-1)$, un petit mot (« où » ?)}

Une telle structure sur \uncertain{$G_{\mathbb{Q}}$} en définit une
analogue sur $G_{\mathbb{C}}$, et réciproquement ; donc une telle
\uncertain{donnée} de structure peut être considérée comme équivalente
\struck{\ill{} partie de \ill{} $X^{+} \subset K(G)$ \ill{}} à la donnée
d'un \struck{cône} cône $K^{+}(G_{\mathbb{C}})$ \struck{$K^{+}(G)$}
\struck{l'ensemble \ill{} contenu de $G_{\mathbb{C}}$}
\struck{($X^{+} + X^{+} \subset X^{+}$)} \struck{stable par produit, tel
que pour tout} dans l'anneau des car.\ virt./$\mathbb{Q}$ $K(G)$, contenant
$0$, \struck{$E$} stable par multiplication, et
\marginal{engendré par les éléments de la base canonique (classes de rep.\
irréd.\ sur $\mathbb{Q}$) qu'il contient.}
À priori, \uncertain{ça} semble ici une donnée de structure faisant
intervenir la structure sur $\mathbb{Q}$. Il y a lieu de prévoir d'autres
axiomes. Par exemple, \ill{} on voudra que pour \struck{$H$}
\struck{\ill{}} $E \in \mathrm{Ob}\,\mathcal{M}^{+}$, $E$ soit à graduations
\emph{positives}.
\note{« l'anneau des car.\ virt./$\mathbb{Q}$ » est interlinéaire ; le
passage biffé mêle plusieurs rédactions successives, raturées l'une sur
l'autre ; la note marginale « engendré par… » est reliée par un trait au
« et » qui la précède}

\textbf{Axiome.} Pour tout \struck{\ill{}} \uncertain{entier} $i$, il n'y a
qu'un nombre fini de types de motifs effectifs de poids $i$ \ill{} (?!).
En fait \ill{} motifs effectifs \ill{} tels que $G^{\circ}$ \ill{}.
\note{encadré à droite, d'une écriture plus petite et plus rapide ; la fin
n'est pas lue}

\marginal{De plus on voudra que la structure envisagée provienne d'une
structure de \uncertain{même} type donnée sur $G^{\circ}$, \struck{\ill{}}
\ill{}.}
\note{note marginale oblique, en bas à gauche ; « même » est écrit comme un
signe $=$}

Un \ill{} $G_{m,\mathbb{C}} \xrightarrow{\ i\ } G_{\mathbb{C}}$
\uncertain{sera dit} \struck{positif} effectif si pour tout
$E \in \mathrm{Ob}\,\mathcal{M}^{+}$, la graduation correspondante de
$E_{\mathbb{C}}$ est à

\page{4}
degrés positifs. On voit que si $i$ est effectif, il en est de même de ses
conjugués par éléments de $G(\mathbb{C})$, et de son conjugué par
l'automorphisme \ill{} de $\mathbb{C}$, \ill{} $\bar{i}$. D'autre part, pour
$i$ donné, la sous-catégorie pleine $\mathrm{Pl}(i)$ de $\mathcal{M}$ formée
des $E$ tels que $E_{\mathbb{C}}$ soit à $i$-graduation positive est
strict.\ pleine, stable par $\otimes$, sommes, par \ill{}, sous-objets, par
$\otimes$, contient $0$, et $\Lambda(-1)$ si on suppose que
\struck{$\varepsilon i(\lambda) = \lambda^{p}$ avec $p$}
$\langle i, \varepsilon \rangle \geq 0$, et vérifie la condition que
$\forall E \in \mathrm{Ob}\,\mathcal{M}$, $\exists n$ tel que
$E(n) \in \mathrm{Ob}\,\mathrm{Pl}(i)$, si on suppose
\[
  \boxed{\langle i, \varepsilon \rangle > 0}
\]
[N.B. on \uncertain{désigne} par $\langle i, \varepsilon \rangle$ l'entier
$\nu$ tel que $\varepsilon(i(\lambda)) = \lambda^{\nu}$].
\note{au-dessus des deux $\langle i, \varepsilon \rangle$ de ce passage, de
sa main : « $\varepsilon(i)$ » ; le premier $\langle i, \varepsilon \rangle$
est en partie surchargé}

Prenant des intersections \struck{si le groupe \ill{} engendré par \ill{}
et les conjugués} \struck{$\mathrm{int}(g) \circ i$ \ill{} les conjugués
$T_{i}$ ($g \in G(\mathbb{C})$), $i$ \ill{}} de sous-catégories
$\mathrm{P}(i)$, on trouve des $\mathcal{M}^{+}$ satisfaisant aux axiomes
considérés, dont peut être \ill{} \ill{} sur la graduation \ill{}.
\marginal{Par intersection de \uncertain{tels} foncteurs, etc.…}
\note{la dernière phrase est interlinéaire}

Les $i$ possibles, $\mathcal{M}^{+}$ étant donné,

\page{6}
\ill{} $\mathcal{M}^{+} \subset \mathrm{Pl}(i)$, \struck{forment}
correspondent à certaines classes de conjugaison d'\uncertain{hom}.

Si $T$ est un tore maximal de $G$, et
$\Gamma = \mathrm{Hom}(G_{m,\mathbb{C}}, T_{\mathbb{C}})$,
$W = (\mathrm{Norm}(T)/T)(\mathbb{C})$,
$W' = (\mathrm{Aut}(G, T)/T)(\mathbb{C})$, on sait que $W \subset W'$,
$W'/W = \mathrm{Autext}(G)(\mathbb{C})$, $W'$ opère sur $\Gamma$, et on
obtient un \emph{cône} \struck{\ill{}} $\Gamma^{+}$ de $\Gamma$, invariant
par $W$ (mais peut-être pas par $W'$ !), contenant $0$, contenu dans
\uncertain{l'hyperplan} $\varepsilon \geq 0$, \struck{un cône itou} un cône
correspondant dans le « quotient » $\Gamma' = \Gamma/W$. Or
\[
  \Gamma' = \Gamma/W
  = \bigl[\mathrm{Hom}(G_{m}, T)/(\mathrm{Norm}_{G}(T)/T)\bigr](\mathbb{C})
  = \bigl[\mathrm{Hom}(G_{m}, G)/G\bigr](\mathbb{C})
\]
\uncertain{est formé des} \uncertain{pts} \struck{à val.\ ds}
$\mathbb{C}$ d'un \uncertain{schéma} \struck{\ill{}} $S$ \ill{} sur
$\mathbb{Q}$, et à ce titre le groupe fondamental
$\pi_{1}(\mathbb{Q}, \mathbb{C})$ opère sur \struck{\ill{}}
$\Gamma = S(\mathbb{C}) = S(\bar{\mathbb{Q}})$, d'ailleurs via un quotient
fini de $\pi_{1}$. Ceci posé, le cône envisagé \struck{contient} est stable
par $\pi_{1}$.
\note{$\mathrm{Aut}$, $\mathrm{Autext}$ et $\mathrm{Hom}$ sont soulignés par
lui (foncteurs en groupes) ; au-dessus de la première égalité, les mots
« une partie » ; dans la dernière égalité il écrit $\Gamma$, sans prime}

\page{8}
\struck{À tout « cône » stable}

Si nous nous intéressons aux structures d'effectivité qui peuvent se
définir \struck{par} en termes d'\struck{\ill{}} une famille de $i$ (comme
intersection des $\mathrm{P}(i)$), \struck{\ill{}} la donnée d'une telle
structure équivaut à celle d'un « cône » dans $\Gamma'$, invariant par
$\pi_{1}$, contenu « strictement » dans celui défini par $\varepsilon$,
contenant la \uncertain{classe} de $i$, satisfaisant à une condition de
\struck{\ill{}} saturation que nous ne préciserons pas ici (condition
automatiquement satisfaite ?).
\marginal{!!! i.e.\ $\varepsilon$ est intérieur au cône polaire !}
\note{« qui peuvent se » et « strictement » sont interlinéaires ; la note
marginale est en face de « strictement »}

Le cas le plus intéressant est celui où $\mathcal{M}^{+}$ \struck{$\otimes$}
est défini comme $\mathrm{Pl}(i)$, où $i$ est comme dans 3°.
\struck{(\ill{})} [N.B., le cône \emph{invariant} engendré contient bien
$j$, et est contenu dans le ½-espace défini par $\varepsilon$, vu \ill{},
$\varepsilon i_{1} = \mathrm{id}$]. On trouve ainsi une application
\[
  I(G) \longrightarrow \mathrm{Eff}(G)
\]
\note{devant $\mathrm{Eff}(G)$, quelques lettres biffées ; « N.B. » est
écrit au-dessus du passage biffé}

\page{10}
\struck{notant que} en particulier, tenant compte de 3°,
\[
  I^{*}(G) \longrightarrow \mathrm{Pol}(G) \times \mathrm{Eff}(G)
\]
\marginal{les $i : G_{m,\mathbb{C}} \to G_{\mathbb{C}}$ définissant une
polarisation, modulo action de $G(\mathbb{R})$…}
\note{cette glose de $I^{*}(G)$ est interlinéaire, sous la formule}

\emph{Questions.} Cette dernière application, ou
$I(G) \to \mathrm{Eff}(G)$, injective ?

Les structures $\in \mathrm{Eff}(G)$ provenant de la théorie des motifs
peuvent-elles toujours se décrire par un $i$, i.e.\ sont-elles dans
$\mathrm{Im}\, I(G)$ ?

(Regarder \ill{} en tordant par un torseur). Examiner aussi, pour ces
questions et celles posées dans 3°, le cas où $G$ est \ill{} de structures
provenant d'un motif sur \ill{} \uncertain{associé} à un
$\xi \in S(\mathbb{Q})$, en utilisant le foncteur de De Rham… ?]
\note{un trait vertical relie en marge ce paragraphe à la note N.B.
ci-dessous}

\marginal{N.B. Le fait que la structure d'effectivité puisse se décrire par
un élément de $I(G)$ (ou $I^{*}(G)$) est une propriété intrinsèque de la
catégorie de motifs $\mathcal{M}$, indép.\ des foncteurs
\uncertain{qu'on utilise}.}
\note{note marginale à gauche, précédée de deux points d'exclamation}

Regarder $G$ commutatif ! \uncertain{$G = \mathrm{Gal}(\bar{\mathbb{Q}},
\mathbb{Q})$} !

Des propriétés arithmétiques plus profondes pourraient être sans doute
formulées par l'introduction de l'idée de \uncertain{filtrations} et
l'interprétation des structures 3° et 4° en termes de \ill{}
\uncertain{ces derniers}.

\page{12}
\struck{Considérons un point $\xi \in S(k)$, ($k$ corps de car.\ $0$), il
définit un foncteur fibre} \struck{$H^{**}_{\xi}$ (Hodge) et
$H^{*}_{\xi}$ (De Rham) : à valeurs dans la catégorie des
\uncertain{vectoriels} sur $k$.}
\note{ces trois lignes sont barrées de grands traits en V ; « $k$ corps de
car.\ $0$ » est interlinéaire}

\emph{Cas d'un tore $T$ sur $\mathbb{Q}$.}
\note{« tore » est souligné}

Défini par un $\pi$-module admissible, libre de t.f.\ sur $\mathbb{Z}$,
soit $M$, ($\pi = \pi_{1}(\mathbb{Q}, \mathbb{C})$).

On aura $\varepsilon \in M^{\pi}$, $i \in \check{M}^{\pi}$,
$\langle i, \varepsilon \rangle = 2$.

Les modules sur $T$ correspondent aux éléments invariants par $\pi$ de
$\mathbb{Z}(M)^{+}$ (combinaisons lin.\ à coeff.\ $\geq 0$).

\struck{Donnons nous $\varepsilon$ \ill{}} La donnée d'une structure de
\uncertain{polarisation} \struck{\ill{}} ne dépend que du tore réel
$T_{\mathbb{R}}$. Ce dernier est de la forme
\struck{$T_{\mathbb{R}} \simeq U \times S$} ext.\ d'un $S$ split par $U$
compact, i.e.\ $U(\mathbb{R})$ compact, \struck{$S$ split}, et
$\varepsilon$ est nécessairement nul sur $U$. Donc on est ramené à $S$, et
il s'agit des structures $\varepsilon_{S}$ et \uncertain{$i_{S}$} sur $S$,
et évidemment $S$ hérite de structure de positivité. J'affirme qu'on doit
avoir $\dim S = 1$, en d'autres termes, que $\operatorname{Ker}
\varepsilon$ est compact.
\note{« polarisation », « ext.\ d'un $S$ split par » et « i.e.\
$U(\mathbb{R})$ compact » sont interlinéaires ; « J'affirme qu'on doit
avoir » est ajouté sous la ligne, devant « $\dim S = 1$ »}

\struck{En effet, si on avait $\dim S \geq 2$, on pourrait \ill{}
(éliminant par \ill{} un \ill{} de $\operatorname{Ker} \varepsilon$)}
\struck{trouver \ill{} \ill{} de dimension \ill{} \ill{}}
\struck{$S = G_{m,\mathbb{R}} \times G_{m,\mathbb{R}}$, avec
$i(\lambda) = (\lambda^{\beta}, \lambda^{\alpha})$,
$\varepsilon(\lambda, \mu) = \mu^{\beta}$,} \struck{$\alpha, \beta$
entiers $\geq 1$, $\alpha\beta = 2$ (donc $\alpha = 1$, $\beta = 2$, ou
$\alpha = 2$, $\beta = 1$).}
\note{tout ce passage est barré de deux grands traits obliques ; lecture
partielle. Les exposants de $i(\lambda)$ et de $\varepsilon(\lambda,\mu)$
sont douteux ; dans la parenthèse finale, « $\alpha = 1$ » est lui-même
biffé}

\page{14}
\emph{Détermination du groupe des \uncertain{relations}
\uncertain{rationnelles} sur $\mathbb{Q}$}

$T_{\mathbb{C}} = D_{\mathbb{C}}(M)$, $M$ réseau où $\pi_{1}$ opère
($\pi_{1} = \pi_{1}(\mathbb{Q}, \mathbb{C})$), via un groupe quotient fini
$\pi$ opérant fidèlement.

$\mathbb{Z}(M)$ anneau des caractères virtuels de $T_{\mathbb{C}}$, $\pi$ y
opère.

$K(T) = \mathbb{Z}(M)^{\pi}$ Anneau de type fini sur $\mathbb{Z}$.

[N.B. $\mathbb{Z}(M)$ est l'anneau affine de $S = D_{\mathbb{Z}}(M)$, sur
lequel $\pi$ opère par automorphismes de groupe. $K(T)$ est l'anneau affine
de $D_{\mathbb{Z}}(M)/\pi$]. Les caractères \emph{effectifs} sont ceux qui
prennent des valeurs $\geq 0$ sur les points de $S(\mathbb{R})^{\circ}$ ??
Comment définir une structure d'effectivité \uncertain{intrinsèque} ?
\marginal{$\lambda_{1}^{u_{1}} \lambda_{2}^{u_{2}} \cdots
\lambda_{r}^{u_{r}}$}
\note{« via un groupe quotient fini $\pi$ opérant fidèlement » et
« virtuels » sont interlinéaires ; la formule marginale est à gauche, en
face de « définir » ; devant $D_{\mathbb{Z}}(M)$, un $\mathbb{Z}$ biffé}

Le procédé standard est de se fixer un $i_{1}^{*} : M \to \mathbb{Z}$,
i.e.\ un $i_{1} : G_{m,\mathbb{C}} \to T_{\mathbb{C}}$ \ill{} les propriétés
habituelles, et donnant donc
\[
  \mathbb{Z}(M) \to \mathbb{Z}(\mathbb{Z})
  \quad \text{d'où} \quad
  G_{m,\mathbb{Z}} \xrightarrow{\ \alpha\ } S
\]
(pas compatible avec $\pi$….). La condition est, \struck{\ill{}} pour un
$\chi \in K(T) \simeq \mathbb{Z}(S)^{\pi}$, que $\chi \circ \alpha \simeq$
\ill{} anneau affine de $G_{m,\mathbb{Z}}$ soit une fonction
\emph{polynômiale}, i.e.\ restant finie sur
$G_{m,\mathbb{Z}}(\mathbb{R}) = \mathbb{R}^{*}$ pour $t \to 0$.

Considérons les caractères de motifs effectifs de poids donné $\rho$,
\ill{} \uncertain{possible} \ill{} \uncertain{qu'il n'y en a qu'un nombre
fini}. Comme \struck{\ill{}} \struck{\ill{}} \ill{} \ill{} quand on passe
de $\mathbb{Q}$ par $\mathbb{R}$, \struck{Alors} on \ill{} si $T$ est de
dim $\geq 3$, la \uncertain{condition} de \uncertain{finitude}
correspondante est \emph{fausse}.
\note{« fausse » est souligné ; la fin du feuillet est serrée et la
lecture du dernier paragraphe est très incertaine}

\section{Équivalence d'Albanese et de Picard des cycles algébriques}

\note{titre de sa main, en haut à droite d'un feuillet de couverture
(feuillet 16) laissé par ailleurs blanc ; le feuillet 17 est blanc}

\page{18}
\textbf{E}
\note{une grande lettre « E », d'une encre plus forte, en tête du feuillet}

$A$ courbe elliptique.

$Y = A \times P$, $P$ courbe non sing.\ complète conn.
\note{l'exposant de $P$ est surchargé (un « 1 » effacé ?), de même dans
$\Theta_{0} \times P$ ci-dessous, où un petit $p$ est écrit sous le $P$}

$\Theta_{0} = \Delta - A \times (0) - (0) \times A$ \struck{\ill{}} diviseur
\uncertain{correspondant} sur $A \times A$ définissant l'autodualité.

$\Theta = \Theta_{0} \times P$ diviseur \struck{correspondant} sur
$A \times A \times P = A \times Y$.

$c$, classe de diviseur sur $Y$, défini par un \struck{fibré} faisceau
inv.\ $L$ comme ($c = \mathrm{cl}(L)$).

$X = \overline{V(L)}$ (clôture projective). $Y^{2} \xrightarrow{\ i\ }
X^{3}$ section nulle.

Le cycle $\Theta$ sur $A \times Y$ définit
\[
  z = (\mathrm{id}_{A} \times i)_{*}(\Theta)
\]
cycle de codim $2$ sur $A \times X^{3}$. On va calculer, si
$s : A \times X^{3} \to A$ est la projection, $s_{*}(z^{2})$, et on va
\uncertain{démontrer} \ill{} des deux \ill{} \ill{} $c$ \ill{} $> 0$,
\ill{} $< 0$.
\note{« $c = \mathrm{cl}(L)$ » est interlinéaire ; les exposants de $Y^{2}$
et $X^{3}$ marquent les dimensions ; la fin de la phrase sur $s_{*}(z^{2})$
n'est pas lue}

Pour cela, notons que si $i' = \mathrm{id}_{A} \times i$,
\[
  z^{2} = i'_{*}(\Theta)^{2} = i'_{*}(\Theta^{2} c')
\]
où $c' = i'^{*} i'_{*}(1) = c^{1}$ (faisceau normal par $i'$)
$= 1_{A} \times c$, $c$ étant le faisceau normal par $i : Y \to X$.

\begin{tikzcd}
  A \times Y \arrow[dr, "t"'] & A \times X \arrow[l, "i'"'] \arrow[d, "s"] \\
  & A
\end{tikzcd}

\note{diagramme en marge gauche ; la flèche $i'$ est dessinée vers
$A \times Y$, bien que le texte pose $i' = \mathrm{id}_{A} \times i$ de
$A \times Y$ dans $A \times X$}

\page{20}
D'où $s_{*}(z^{2}) = t_{*}(\Theta^{2} c')$, où $t : A \times Y \to A$ est la
projection. \ill{} \ill{}, \ill{} \ill{} :
\[
  K(z^{2}) = K(\Theta^{2} c') = K(\Theta^{2} \cdot 1_{A} \times c)
  = K(p_{*}(\Theta^{2}) \cdot c)
\]
où $p$ est la projection $A \times Y \to Y$.
\note{la lettre notée $K$ est lue sous réserve ; devant $K(z^{2})$ une
lettre surchargée ; en marge gauche, la flèche verticale
$t : A \times Y \to A$}

Nous allons calculer $p_{*}(\Theta^{2})$, en remarquant que
$p = p' \times \mathrm{id}_{P}$, où $p' : A \times A \to A$ est la deuxième
projection.

\begin{tikzcd}[row sep=small]
  A \times Y \arrow[r, "p"] & Y \\
  \parallel & \\
  A \times A \times P \arrow[r] & A \times P
\end{tikzcd}

\[
  \Theta = \Theta_{0} \times 1_{P} \quad \text{d'où} \quad
  \Theta^{2} = \Theta_{0}^{2} \times 1_{P} \quad \text{d'où}
\]
\[
  p_{*}(\Theta^{2}) = p'_{*}(\Theta_{0}^{2}) \times
  \mathrm{id}_{P*}(1_{P}) = p'_{*}(\Theta_{0}^{2}) \times 1_{P} .
\]
D'autre part on a
\[
  \Theta_{0}^{2} = \Delta^{2} + 0 + 0 - 2(0,0) - 2(0,0) + 2(0,0)
  = -2(0,0)
\]
compte tenu que $\Delta^{2} = 0$ (équiv.\ alg.) puisque $\chi(A) = 0$.
\note{« (équiv.\ alg.) » est interlinéaire}

D'où
\[
  p_{*}(\Theta^{2}) = -2(0 \times 0 \times P)
\]
et on trouve finalement
\[
  K(s_{*}(z^{2})) = -2\, c(0,0)
\]
si on pose $c(0,0) = K(c \cdot (0 \times 0 \times P))$.

Il existe un $c$ tel que $c(0,0) = 1$, p.\ ex.\
$c = A \times A \times (a)$, où $a \in P(k)$. \emph{Donc} remplaçant $c$
par $-c$, on trouve comme classes \uncertain{numériquement} $-2$ et $2$, qui
ont des signes différents !!

\end{document}
